\clearpage
\section{Schur reduction and coercivity of the Weil form}
\label{sec:schur-coercividad}
% Fuente íntegra: ampliacion/expediente/AGENTE_SCHUR_DETALLES.md

% PROCEDENCIA 11 de septiembre de 2026. Desarrollo focal para integración por el editor único.
% PROCEDENCIA No modifica los manuscritos I–VIII ni sus PDF.

\subsection{Basis, domain, and result obtained}
\label{schur-add-sub-6}

We use the full columns already constructed from the native arithmetic
of APP, the oriented TRIT state, the TPK operations and memory, and the
joint structure of the continuum. The native product expresses irreducibles and
powers; their lengths and weights subsequently feed the prime,
gamma, and polar readouts. The internal coordinate of \(\pi\) and the rational readout of the
Euler--Mascheroni finite part enter after that representation. No
zero, desired value of the functional, or metrological constant selects these
operations.

We retain the full columns and the domains constructed in the
chapter on the double circle and spectral transport, together with the
support and detail estimates. The columns and the connector are also made explicit in
Section~\ref{sec:lectores-gamma}. The joint construction
of the continuum remains the common antecedent: it is not replaced by this
functional decomposition, nor are its five constructions replaced by five modules.

The new operation assembled here is the exact elimination of the
infinite detail space by means of its coercive form. It produces a finite matrix
with all coupling corrections, not merely a matrix of
means. A specific case is also settled:

For every function in the closure of
\(C_c^\infty((-1/18,1/18))\) in the joint norm of the readouts,
\begin{equation}
\mathscr W(f,f)\geq\frac12\|f\|_2^2.
\label{schur-add-display-1}
\end{equation}
The same assertion holds on any translated interval of length
\(1/9\), and on all its interior nonadic refinements.

The full window has length \(1/9>0.085626\ldots\), the endpoint
of the preceding baseline criterion. It contains no active prime translation,
since \(1/9<\log2\). This limitation of the example remains explicit.

\textbf{Provenance.} Authorial architecture and readouts: pre-existing. Detail
coercivity: a result recovered from the zero-mean threshold theorem, retaining an
intermediate estimate stronger than its abbreviated logarithmic bound. Schur assembly,
residual certification, and proof for the interval \(1/9\): formalization and
certificate added in this section; no global historical novelty is claimed.

\subsection{Partition, cell means, and details}
\label{schur-add-sub-46}

Fix \(R>0\), \(I_R=(-R/2,R/2)\), and \(H_R=L^2(I_R)\), with extension
by zero. Let \(V_R\) be the closure of \(\mathcal D_R=C_c^\infty(I_R)\) in the
joint norm of the full readouts. Equivalently, its norm may be taken to be
\begin{equation}
\begin{aligned}
\|f\|_{V_R}^2&=\|f\|_2^2+Q_\Gamma(f),\\
Q_\Gamma(f)&=\int_0^\infty\mu(a)\|(I-T_a)f\|_2^2\,da,\\
\mu(a)&=\frac{e^{-a/2}}{1-e^{-2a}}.
\end{aligned}
\label{schur-add-display-2}
\end{equation}
The prime and polar channels are bounded on each fixed window.
The gamma form is closed: it is the form of a positive Fourier
multiplier, restricted to the subspace of functions supported in the
window; equivalently, this is verified by convergence in the
direct integral of differences. The full form is a bounded perturbation
of it and defines a semibounded self-adjoint operator \(L_R\).

Partition \(I_R\) into contiguous intervals \(I_j=(\ell_j,r_j)\), with
lengths \(h_j\leq h\), and set
\begin{equation}
u_j=h_j^{-1/2}\mathbf1_{I_j},\qquad
Ua=\sum_{j=1}^M a_j u_j,\qquad P=UU^*,\quad Q=I-P.
\label{schur-add-display-3}
\end{equation}
In the last sum, \(a=(a_1,\ldots,a_M)\) is the coefficient vector;
the endpoints \(\ell_j,r_j\) are used only in defining \(I_j\).

The indicator functions belong to \(V_R\): their gamma energy is controlled by
\(\|(I-T_a)u_j\|_2^2=2\min(|a|/h_j,1)\). Thus \(P\) acts
continuously on \(V_R\), and every \(f\in V_R\) has the unique decomposition
\begin{equation}
f=Ua+d,\qquad a_j=\langle f,u_j\rangle,\qquad
d\in V_D:=V_R\cap QH_R.
\label{schur-add-display-4}
\end{equation}
Here \(d\) has zero integral on each cell. The space \(V_D\) is
closed in the form norm and dense in \(QH_R\).

The estimate for zero-mean details gives a bound \(\mathscr W_R(d,d)\geq\eta\|d\|_2^2\)
if a sufficiently fine partition is chosen. One may use its
expressed logarithmic threshold or retain the exact intermediate estimate
\begin{equation}
\begin{aligned}
\eta_N(h,R)={}&4\sum_{n=0}^N\frac1{4n+1}
-\frac{h^2(N+1)(2N+1)}{\pi^2}\\
&+c_\Gamma-2\sum_{k\log p\leq R}(\log p)p^{-k/2}-\beta_R,\\
\beta_R={}&2\sinh(R/2)-R.
\end{aligned}
\label{schur-add-eq-S1}
\end{equation}
\begin{lemma}[Retained bound for cell details]
\label{schur-add-retencion-detalles}
For \(d\in V_D\) and every integer \(N\geq0\),
\[
\mathscr W_R(d,d)\geq\eta_N(h,R)\|d\|_2^2.
\]
\end{lemma}
\begin{proof}
On each cell \([\ell_j,r_j]\), define
\(F(x)=\int_{\ell_j}^x d(t)\,dt\). The zero mean gives
\(F(\ell_j)=F(r_j)=0\). The assembly of these primitives, extended
by zero, belongs to \(H_0^1(I_R)\), satisfies \(F'=d\), and, by
the Dirichlet inequality on each cell,
\begin{equation}
\|F\|_2^2\leq\frac{h^2}{\pi^2}\|d\|_2^2.
\label{schur-add-primitiva}
\end{equation}
The positive expansion
\(\mu(a)=\sum_{n\geq0}e^{-\lambda_n a}\), \(\lambda_n=2n+1/2\),
allows us to write \(Q_\Gamma=\sum E_{\lambda_n}\), where
\begin{equation}
E_\lambda(d)=\int_0^\infty e^{-\lambda a}
 \|(I-T_a)d\|_2^2\,da
 =\frac2\lambda\|d\|_2^2-\langle d,C_\lambda d\rangle\geq0,
\label{schur-add-energia-lambda}
\end{equation}
and \(C_\lambda\) is convolution with \(e^{-\lambda|x|}\).
For the unitary transform
\(\mathcal F_u=(2\pi)^{-1/2}\widehat{\phantom f}\), the multiplier
of \(C_\lambda\) is \(2\lambda/(\lambda^2+\xi^2)\). By Plancherel,
\begin{equation}
\begin{aligned}
\langle d,C_\lambda d\rangle
&=\int_{\mathbb R}\frac{2\lambda\xi^2}{\lambda^2+\xi^2}
 |\mathcal F_uF(\xi)|^2\,d\xi\\
&\leq2\lambda\|F\|_2^2
 \leq\frac{2\lambda h^2}{\pi^2}\|d\|_2^2 .
\end{aligned}
\label{schur-add-cota-convolucion}
\end{equation}
Retaining \(N+1\) summands and discarding only the nonnegative tail gives
\[
Q_\Gamma(d)\geq
\left[4\sum_{n=0}^N\frac1{4n+1}
-\frac{h^2(N+1)(2N+1)}{\pi^2}\right]\|d\|_2^2.
\]
The scalar term contributes \(c_\Gamma\|d\|_2^2\).
Each prime translation contributes at least
\(-2w_{p,k}\|d\|_2^2\), by Cauchy--Schwarz and the unitarity
of \(T_{k\log p}\). Only lengths less than or equal to
\(R\) enter.
Finally, the polar contribution is
\[
2\left|\int_{I_R}d(x)\cosh(x/2)\,dx\right|^2
-2\left|\int_{I_R}d(x)\sinh(x/2)\,dx\right|^2
\geq-\beta_R\|d\|_2^2,
\]
since \(2\int_{I_R}\sinh^2(x/2)\,dx=2\sinh(R/2)-R\).
Summation proves the bound.
\end{proof}
There is no need to replace the harmonic sum by its lower integral bound.
For the certificate on the window \(1/9\), one also checks
\((2N+1/2)h<\pi\), so that each retained gamma bound is nonnegative.

For any fixed \(R\), there is a partition that makes this
bound positive. Indeed, if \(0<h\leq1\) and \(N=\lfloor1/h\rfloor\), then
\[
4\sum_{n=0}^N\frac1{4n+1}\geq\log(4N+5)\geq\log(1+4/h),
\qquad
\frac{h^2(N+1)(2N+1)}{\pi^2}<\frac23.
\]
The first term diverges as \(h\downarrow0\); the prime,
scalar, and polar terms remain fixed by \(R\). This preserves
the distinction between coercivity of the details and the sign of the
full form.

\subsection{The exact Schur complement has finite size}
\label{schur-add-sub-97}

Let \(D\) be the operator associated with the restriction of the form to \(V_D\).
If \(\eta>0\), then \(D\geq\eta I\) on \(QH_R\),
\(D^{-1}\) exists, and \(\|D^{-1}\|\leq\eta^{-1}\).
We do not presume that \(L_R\) preserves the detail space.

The indicator functions also belong to the operator domain of \(L_R\).
The expression for their gamma image is given in the next section: it has only
logarithmic singularities, which belong to \(L^2(I_R)\).
This allows us to define, without improperly extending a form functional,
\begin{equation}
A=U^*L_RU,\qquad B=QL_RU:\mathbb C^M\longrightarrow QH_R.
\label{schur-add-display-6}
\end{equation}
The inner product is linear in the first variable, as in the
manuscript. The matrices are written so that
\(\mathscr W_R(Ua,Ua)=a^*Aa\). The block identity is
\begin{equation}
\mathscr W_R(Ua+d,Ua+d)
=a^*Aa+2\operatorname{Re}\langle Ba,d\rangle
+\mathfrak d[d,d],
\label{schur-add-display-7}
\end{equation}
where \(\mathfrak d\) is the closed form of \(D\).

\begin{theorem}[Exact elimination of the detail space]
\label{schur-add-schur-exacto}
The matrix
\begin{equation}
S=A-B^*D^{-1}B
\label{schur-add-eq-S2}
\end{equation}
satisfies the exact identity
\begin{equation}
\boxed{\mathscr W_R(Ua+d,Ua+d)
=a^*Sa+
\mathfrak d[d+D^{-1}Ba,d+D^{-1}Ba].}
\label{schur-add-eq-S3}
\end{equation}
In particular, \(\mathscr W_R\geq0\) on \(V_R\) if and only if
\(S\succeq0\). Its negative index equals that of \(S\).
\end{theorem}
\begin{proof} For \(z=D^{-1}Ba\in\operatorname{Dom}D\),
\(\mathfrak d[z,d]=\langle Ba,d\rangle\). Expanding the second term
of \eqref{schur-add-eq-S3} and using \(\mathfrak d[z,z]=a^*B^*D^{-1}Ba\) gives the equality.
Sufficiency is immediate. For necessity, choose
\(d=-D^{-1}Ba\). The triangular change \((a,d)\mapsto(a,d+D^{-1}Ba)\)
is invertible on \(\mathbb C^M\oplus V_D\). The resulting second form
is positive definite; the negative directions are exactly those of
the finite matrix.
\end{proof}

This theorem converts the finite-negative-index result into an exact matrix. The matrix still contains \(D^{-1}\): finite size does not
mean that its entries have already been calculated. The next section
provides a way to enclose them with controlled residuals.

\subsection{Explicit images and certifiable bounds for the corrections}
\label{schur-add-sub-148}

Set \(M(x)=\int_x^\infty\mu(a)\,da\), for \(x>0\).
For an indicator function \(u=h^{-1/2}\mathbf1_{(a,b)}\), its gamma image,
restricted to \(I_R\), is almost everywhere
\begin{equation}
(L_\Gamma u)(x)=h^{-1/2}
\begin{cases}
M(x-a)+M(b-x),&a<x<b,\\
-M(a-x)+M(b-x),&x<a,\\
-M(x-b)+M(x-a),&x>b.
\end{cases}
\label{schur-add-eq-S4}
\end{equation}
The equality is obtained by integrating the two translations: inside the
cell there remains one tail for each endpoint, and outside there remains the integral over the
unique interval of translations that meets the cell. It can first be proved
with the regulator \(a_{\rm traslado}\geq\epsilon\); convergence in
\(L^2(I_R)\) of the logarithmic expression justifies its weak form
limit. The other parts are
\begin{equation}
L_{\rm primo}u=-\sum_{k\log p\leq R}w_{p,k}
M_R(T_{k\log p}+T_{-k\log p})u,
\label{schur-add-display-11}
\end{equation}
\begin{equation}
\begin{aligned}
L_{\rm polar}u(x)&=2\int_a^b h^{-1/2}\cosh((x-y)/2)\,dy,\\
L_Ru&=L_\Gamma u+c_\Gamma u+L_{\rm primo}u+L_{\rm polar}u.
\end{aligned}
\label{schur-add-eq-S5}
\end{equation}
All discontinuities and singularities of \eqref{schur-add-eq-S4}--\eqref{schur-add-eq-S5} lie in a
finite list of endpoints and endpoints translated by expressed clocks.
Thus, \(A_{ij}\) and
\begin{equation}
G_B=B^*B,\qquad (G_B)_{ij}=\langle Be_j,Be_i\rangle
\label{schur-add-display-13}
\end{equation}
are calculated by explicit one-dimensional integrals. In this last
formula, \(e_j\) denotes the standard coordinate vector of
\(\mathbb C^M\), not the indicator function in \(H_R\).
The initial bound is
\begin{equation}
A-\eta^{-1}G_B\preceq S\preceq A.
\label{schur-add-eq-S6}
\end{equation}
This is an operator inequality with the cross terms preserved, not a
replacement of \(G_B\) by its diagonal entries.

\begin{theorem}[Residual certification of the complement]
\label{schur-add-residual}
For any finite
approximation operator \(Z:\mathbb C^M\to\operatorname{Dom}D\), define
\begin{equation}
E=B-DZ,\qquad C_Z=Z^*B+B^*Z-Z^*DZ.
\label{schur-add-display-15}
\end{equation}
Then
\begin{equation}
B^*D^{-1}B=C_Z+E^*D^{-1}E,
\label{schur-add-display-16}
\end{equation}
and therefore
\begin{equation}
\boxed{A-C_Z-\eta^{-1}E^*E\preceq S\preceq A-C_Z.}
\label{schur-add-eq-S7}
\end{equation}
\end{theorem}
\begin{proof}
Substitute \(B=DZ+E\) and expand the product. All
terms are defined because the columns of \(Z\) belong to
\(\operatorname{Dom}D\). The bound uses \(0<D^{-1}\leq\eta^{-1}I\).
\end{proof}

An effective choice of \(Z\) consists in solving the
Galerkin system on step-function details of a nonadic refinement. Its
columns are finite combinations of indicator functions and have images
\eqref{schur-add-eq-S4}--\eqref{schur-add-eq-S5}, so \(E\) and its Gram are explicitly integrable. The
Galerkin system is not declared an exact solution: the term
\(\eta^{-1}E^*E\) retains its error on the infinite complement.

To certify these integrals without losing the singular endpoints,
split
\begin{equation}
M(x)=-\tfrac12\log x+\log2+\tfrac\pi4+r(x),
\qquad r(0)=0,\quad r'(x)=\tfrac1{2x}-\mu(x).
\label{schur-add-eq-S8}
\end{equation}
On \(0<x\leq R\),
\begin{equation}
-\tfrac14 e^{R/2}\leq r'(x)\leq\tfrac14,
\qquad |r'(x)|\leq\tfrac14e^{R/2}.
\label{schur-add-display-19}
\end{equation}
Indeed, \(\mu(x)=e^{x/2}/(2\sinh x)\) and
\(x\leq\sinh x\leq xe^x\) give
\(e^{-x/2}/(2x)\leq\mu(x)\leq e^{x/2}/(2x)\).
The inequalities \(1-e^{-y}\leq y\) and \(e^y-1\leq ye^y\)
prove the bounds. We do not assert that \(r'\) has constant sign
on an arbitrarily large window.
We integrate \(\log x\) and \(\log^2x\) analytically on the small
endpoint segments; the remainder \(r\) is bounded by its Lipschitz constant.
Away from the endpoints, subdivided interval quadratures and
the exponential series for \(M\) are used. For example,
\begin{equation}
\int_0^\delta\log^2x\,dx
=\delta[(\log\delta)^2-2\log\delta+2].
\label{schur-add-display-20}
\end{equation}
Products of two singularities at distinct points are handled
by separating their neighborhoods; for a shared endpoint, the same
quadratic integral is used. The resulting intervals are inserted into \eqref{schur-add-eq-S7}, and
positivity is verified by an \(LDL^*\) decomposition with
strictly positive interval pivots or a perturbation bound
for a rational matrix. This specifies a finite certificate
with control of the infinite complement; no mesh without a tail bound is used.

\subsection{Coercivity on a window of length 1/9}
\label{schur-add-sub-248}

Take \(R=h=1/9\), a single cell, and \(u=R^{-1/2}\mathbf1_{I_R}\).
This is a full interval, not nine intervals separated by gaps. The detail
section is all of \(V_R\cap\{u\}^{\perp}\).

\subsubsection{Margin of the infinite block}
\label{schur-add-sub-254}

In \eqref{schur-add-eq-S1}, take \(N=13\). There are no prime terms because
\(R<\log2\). This gives
\begin{equation}
\eta=4\sum_{n=0}^{13}\frac1{4n+1}
-\frac{(1/9)^2\,14\cdot27}{\pi^2}+c_\Gamma-\beta_{1/9}>1.
\label{schur-add-eq-S9}
\end{equation}
The interval certificate gives
\(1.00363976543713133046<\eta<1.00363976543713133048\).
The condition \((53/2)/9<\pi\) also guarantees the positivity of
each retained gamma bound. The sum is preserved; replacing it by the
bound \(\log(1+4/h)-2/3\) would lose the necessary margin here.

\subsubsection{Mean block}
\label{schur-add-sub-269}

Inside the window, writing \(x=R(t-1/2)\),
\begin{equation}
L_Ru(x)=R^{-1/2}g(t),\quad
g(t)=c_\Gamma+M(Rt)+M(R(1-t))
+8\sinh(R/4)\cosh(R(t-1/2)/2).
\label{schur-add-display-22}
\end{equation}
Thus \(A=\int_0^1g(t)\,dt\), and \(\|B\|^2\) is the variance
of \(g\) on the unit interval. Applying Tonelli to \(M\) gives
\begin{equation}
A=c_\Gamma+\frac2R\sum_{n=0}^\infty
\frac{1-e^{-(2n+1/2)R}}{(2n+1/2)^2}
+\frac{32}R\sinh^2(R/4).
\label{schur-add-eq-S10}
\end{equation}
If the first \(L\) terms are retained, the positive tail is at
most \(\lambda_L^{-2}+(2\lambda_L)^{-1}\), with
\(\lambda_L=2L+1/2\). The first term of the tail is separated, and the
remainder is bounded by the integral of a decreasing function. For
\(L=1024\), the certificate yields
\begin{equation}
0.9723307136651774<A<0.9767284617331572,
\qquad A>97/100.
\label{schur-add-eq-S11}
\end{equation}
A tighter approximation is not needed for the proof.

\subsubsection{Coupling between means and details}
\label{schur-add-sub-295}

On \(0<x\leq1/9\), the estimate in \eqref{schur-add-eq-S8} improves to
\begin{equation}
-9/10\leq r'(x)\leq0.
\label{schur-add-display-25}
\end{equation}
Indeed, for \(0<x\leq1/9\),
\begin{equation}
\frac{\sinh x}{x}
=\sum_{k\geq0}\frac{x^{2k}}{(2k+1)!}
\leq\frac1{1-x^2}\leq1+\frac x2\leq e^{x/2}.
\label{schur-add-display-26}
\end{equation}
The intermediate inequality is equivalent to \(1-2x-x^2\geq0\),
which holds on that interval. Thus \(\mu(x)\geq1/(2x)\).
On the other hand,
\begin{equation}
\mu(x)-\frac1{2x}\leq\frac{e^{x/2}-1}{2x}
\leq\frac14 e^{x/2}
\leq\frac1{4(1-1/18)}=\frac9{34}<\frac9{10}.
\label{schur-add-display-27}
\end{equation}
We retain the less sharp bound \(9/10\) used by the certificate.
The function \(q(t)=r(Rt)+r(R(1-t))\) is symmetric about
\(1/2\), and its Lipschitz constant is at most \(9R/10\).
Its oscillation is at most \(9R/20=1/20\); hence its standard
deviation is at most \(1/40\). This last assertion follows from
\(\operatorname{Var}(X)\leq(\sup X-\inf X)^2/4\): subtract the
midpoint of the range and use the fact that the mean minimizes the squared error.

The singular part \(v(t)=-\tfrac12\log(t(1-t))\) has mean 1 and
variance \(1-\pi^2/12\). To verify this, integrate
\(\log t\), \(\log^2t\), and the positive series for
\(\log t\log(1-t)\), which gives
\(\int_0^1\log t\log(1-t)\,dt=2-\sum_{n\geq1}n^{-2}\).
The identity \(\sum n^{-2}=\pi^2/6\) can be obtained by applying Parseval
to the sine series of \(x\) on \((-\pi,\pi)\): its
coefficients are \(2(-1)^{n+1}/n\), and
\((1/\pi)\int_{-\pi}^{\pi}x^2dx=2\pi^2/3\). This is a subsequent
functional recognition of the same \(\pi\), not a selection of its generator.

The polar part has oscillation
\(8\sinh(R/4)[\cosh(R/4)-1]\), and standard deviation at most
half that value. The triangle inequality in the space of centered
functions gives, for the full coupling,
\begin{equation}
\boxed{\|B\|^2\leq
\left[\sqrt{1-\pi^2/12}+\frac1{40}
+4\sinh(R/4)(\cosh(R/4)-1)\right]^2<\frac15.}
\label{schur-add-eq-S12}
\end{equation}
The scalar term and all constants independent of \(t\)
are removed by centering, not by a supposed disappearance of the
cross terms. The interval evaluation of the bounding expression is
\(0.19926357337268538427\ldots<1/5\).

\subsubsection{Positive Schur complement and coercivity of the entire form}
\label{schur-add-sub-348}

\begin{theorem}[Coercivity on the full nonadic window]
\label{schur-add:coercividad-ventana}
On \(I_{1/9}=(-1/18,1/18)\), every function in the joint domain
\(V_{1/9}\) satisfies
\[
\mathscr W_{1/9}(f,f)\geq\tfrac12\|f\|_2^2.
\]
The bound holds on every translated interval of the same length and on all
its interior refinements.
\end{theorem}
\begin{proof}
From \eqref{schur-add-eq-S6}, \eqref{schur-add-eq-S9},
\eqref{schur-add-eq-S11}, and \eqref{schur-add-eq-S12},
\begin{equation}
S=A-B^*D^{-1}B>\frac{97}{100}-\frac15=\frac{77}{100}>0.
\label{schur-add-display-29}
\end{equation}
Moreover, for \(f=au+d\),
\begin{equation}
\mathscr W_R(f,f)-\tfrac12\|f\|_2^2
\geq(A-\tfrac12)|a|^2
-2\|B\||a|\|d\|_2+(\eta-\tfrac12)\|d\|_2^2.
\label{schur-add-display-30}
\end{equation}
The scalar matrix on the right-hand side has positive diagonal and
determinant strictly greater than
\(\tfrac{47}{100}\cdot\tfrac12-\tfrac15=\tfrac7{200}>0\).
This proves \(\mathscr W_R(f,f)\geq\tfrac12\|f\|_2^2\)
for every \(f\in V_R\). The argument controls all the infinitely many
details; it does not replace them by finitely many indicator functions. Translation
preserves the correlations and multiplies the polar evaluations
by \(e^{t/2}\) and \(e^{-t/2}\), whose factors cancel in the cross products.
It therefore also preserves the full form and the bound.
\end{proof}

\subsection{Refinement, double circle, and the two moments}
\label{schur-add-sub-367}

Each subdivision into nine preserves exactly
\(u_{\rm padre}=\tfrac13\sum_{r=0}^8u_r\), and the recomposition of
functions occurs before applying \(J_+\) and \(J_-\). Therefore,
the Grams of all nonadic indicator functions inside the window
satisfy \(G_n\geq\tfrac12 I\) and
\(V_n^*G_{n+1}V_n=G_n\). The eight details per predecessor and
all their cross terms are preserved. This is uniformity in depth on this window,
not a uniform result as \(R\) increases.

For a general partition, the exact Schur matrix does not automatically
obey the same compression as the raw Gram. If one refines and expresses the
fine Schur complement in old-mean/new-detail coordinates, the old Schur
complement is obtained by also eliminating that new detail. This is
associativity of minimization in \eqref{schur-add-eq-S3}. Asserting
\(S_{\rm padre}=V^*S_{\rm hijos}V\) would omit that second correction.

On the certified window, the connector already constructed,
\(\mathcal C_R=J_{-,R}L_R^{\rm cola}\), where
\(L_R^{\rm cola}\) is the cotrace
\eqref{gamma-add-eq-11} applied to the gamma component, satisfies
\(\mathcal C_RJ_{+,R}=J_{-,R}\). Its restriction to the closure of the
effective image \(\mathcal M_{+,R}\) is contractive, because
\begin{equation}
\|J_{-,R}f\|^2\leq\|J_{+,R}f\|^2.
\label{schur-add-display-31}
\end{equation}
Contractivity is now proved by means of Schur; it is not defined
from a putatively positive Gram. Its defect
\(\Delta_R=(I-\mathcal C_R^*\mathcal C_R)^{1/2}\) is then taken on that
effective space, and we construct
\begin{equation}
\Xi_Rf=J_{-,R}f\oplus\Delta_RJ_{+,R}f.
\label{schur-add-display-32}
\end{equation}
If \(P_-\) is the projection onto the first summand,
\begin{equation}
\Xi_R^*\Xi_R=J_{+,R}^*J_{+,R},\qquad
\Xi_R^*P_-\Xi_R=J_{-,R}^*J_{-,R}.
\label{schur-add-display-33}
\end{equation}
These are equalities of forms on the joint domain, not products of
unbounded operators on all of \(L^2\). They give the two moments in
the dilation codomain. An additional identification with a
specific \(T_N\) of the double circle requires the intertwining
with that operator: an arbitrary projection is not identified with
\(T_N\) merely because the two moments are available.

\subsection{Reproducibility checks and conditions for extension}
\label{schur-add-sub-411}

The verification program retained in
Appendix~\ref{ap:procedencia-reproduccion} certifies the constants with interval arithmetic and explicit
tail bounds. It uses the internal prefix of \(\pi\) with its hashes and the
previously retained rational producer of \(\gamma\) and \(\log2\). It does not use \texttt{\detokenize{mpmath.pi}},
\texttt{\detokenize{mpmath.euler}}, zero tables, or floating-point arguments as inputs.

% DOCUMENTACION ```text
% DOCUMENTACION python3 verificar_schur_ventana_nonadica.py --receipt CERTIFICADO_SCHUR_R_UN_NOVENO.json
% DOCUMENTACION python3 -O verificar_schur_ventana_nonadica.py --receipt CERTIFICADO_SCHUR_R_UN_NOVENO_OPTIMIZADO.json
% DOCUMENTACION ```

Both runs satisfy the same 21 interval checks.
% ESTADO HISTORICO: PASS_SCHUR_FIXED_WINDOW_ONE_NINTH.
The receipts retain exact rational endpoints, precision, library
version, provenance, and hashes. The adversarial check
\(\begin{pmatrix}1&2\\2&1\end{pmatrix}\) has two positive diagonal
blocks and Schur complement \(-3\): it rejects using their signs alone.

Substantive falsifiers: omitting the negative polar term; interchanging cell
means with point sampling; losing a prime translation on a
larger window; replacing the Schur complement by \(A\); accepting Galerkin without
\(E^*E/\eta\); or conflating all interior refinements with
all windows. None is part of this certificate.

The next concrete extension is a window \(R>\log2\), with
a nonadic partition whose detail block is coercive. There,
\eqref{schur-add-eq-S7} must be evaluated, including at least the clock \(\log2\).
The present result asserts neither that new sign, nor global Weil
positivity, nor RH. Nor does it infer a negation of the corpus from
this local scope: it provides an exact elimination, a certifiable
procedure, and a first settled case extending the preceding baseline case.

% DOCUMENTACION ## Localizadores materiales
% DOCUMENTACION 
% DOCUMENTACION - `output/ARTICULO_VIII_PRIMOS_ESTRUCTURA_ESPECTRAL_20260910/manuscrito/03_DOBLE_CIRCULO_Y_TRANSPORTE_ESPECTRAL.md`: §§2–3, líneas 120–260; §12, 1000–1244; §13, 1245–1514; §14, 1515–1739; §15, 1740–1896.
% DOCUMENTACION - Mismo paquete, `fuente/pruebas/python/verificar_gram_nueve_ventanas.py`: publicación autenticada de π, forma intervalar y cola gamma.
% DOCUMENTACION - Mismo paquete, `fuente/pruebas/python/partes_finitas_exactas.py`: productores racionales de γ y logaritmos.
% DOCUMENTACION - Mismo paquete, `fuente/pruebas/python/verificar_nueve_intervalos_coercividad.py`: estimación completa previa en soportes desconectados.
% DOCUMENTACION - Mismo paquete, `gestion/recibo_gamma/RECIBO_CONSTANTES_GAMMA.json` y `RECIBO_GENEALOGIA_GAMMA.json`: procedencia heredada; esta sección no cambia su estado ni pretende estar ya incorporada en esos recibos.
% DOCUMENTACION - Serie activa: integral de 2.249 páginas y desarrollo estructural común. El gate global de primera aplicación sigue devolviendo el testigo de 2.084 páginas; conforme a continuidad, no desplaza la selección autoral de la serie ni este delta focal.
% DOCUMENTACION
